\documentclass[11pt]{article}
\usepackage{amsmath}
\usepackage{amssymb}
\usepackage{geometry}
\usepackage{enumitem}
\usepackage{xcolor}

\geometry{margin=1in}
\setlength{\parskip}{0.5em}
\setlength{\parindent}{0pt}

\title{Mathematical Notation Guide for peppwR Methods}
\author{Explanatory Guide}
\date{\today}

\begin{document}

\maketitle

\section*{Introduction}

This guide explains all mathematical formulas and notation used in the Methods section of the peppwR JSS article. Each formula is presented with plain-language explanations of what the symbols mean, how to read the notation, what the formula does conceptually, and why it's used in the peppwR context.

\section{Statistical Framework Section}

\subsection{Formula 1: Statistical Power Definition}

\textbf{Mathematical notation:}
\[P(\text{reject } H_0 | H_1 \text{ true}) = 1 - \beta\]

\textbf{What the symbols mean:}
\begin{itemize}
\item $P(\cdot)$ = probability of something happening
\item $H_0$ = null hypothesis (no difference between groups)
\item $H_1$ = alternative hypothesis (there is a difference between groups)
\item $\beta$ = Type II error rate (probability of missing a real effect)
\item $|$ = "given that" or "conditional on"
\end{itemize}

\textbf{How to read it:} "Power equals the probability of rejecting the null hypothesis, given that the alternative hypothesis is true, which equals 1 minus beta"

\textbf{What it does:} Defines statistical power as the chance of correctly detecting a real effect when one actually exists

\textbf{Why peppwR uses it:} This is the fundamental concept that peppwR calculates through simulation - the probability that your experiment will successfully detect a real biological effect

\subsection{Formula 2: Hypothesis Testing Framework}

\textbf{Mathematical notation:}
\[H_0: \mu_{\text{treatment}} = \mu_{\text{control}} \quad \text{vs} \quad H_1: \mu_{\text{treatment}} \neq \mu_{\text{control}}\]

\textbf{What the symbols mean:}
\begin{itemize}
\item $\mu_{\text{treatment}}$ = mean abundance in the treatment group
\item $\mu_{\text{control}}$ = mean abundance in the control group
\item $=$ = equals (no difference)
\item $\neq$ = not equals (there is a difference)
\end{itemize}

\textbf{How to read it:} "Null hypothesis: treatment mean equals control mean, versus alternative hypothesis: treatment mean does not equal control mean"

\textbf{What it does:} Sets up the statistical question we're trying to answer - is there a difference between groups?

\textbf{Why peppwR uses it:} This is the standard framework for comparing peptide abundances between experimental conditions

\subsection{Formula 3: Sample Size Determination}

\textbf{Mathematical notation:}
\[n^* = \min\{n : P(\text{reject } H_0 | \delta, n) \geq \pi_{\text{target}}\}\]

\textbf{What the symbols mean:}
\begin{itemize}
\item $n^*$ = optimal sample size
\item $\min\{\cdot\}$ = minimum value that satisfies the condition
\item $n$ = sample size (number of samples per group)
\item $\delta$ = effect size (how big the difference is)
\item $\pi_{\text{target}}$ = target power level (usually 0.8 for 80\% power)
\item $\geq$ = greater than or equal to
\end{itemize}

\textbf{How to read it:} "n-star equals the minimum n such that the probability of rejecting the null hypothesis, given delta and n, is greater than or equal to the target power"

\textbf{What it does:} Finds the smallest sample size that will give you your desired statistical power

\textbf{Why peppwR uses it:} This answers the key question "How many samples do I need?" for your experiment

\subsection{Formula 4: Power Estimation}

\textbf{Mathematical notation:}
\[\pi(n, \delta) = P(\text{reject } H_0 | \delta, n)\]

\textbf{What the symbols mean:}
\begin{itemize}
\item $\pi(n, \delta)$ = power as a function of sample size and effect size
\item $P(\cdot)$ = probability
\item The rest as defined above
\end{itemize}

\textbf{How to read it:} "Pi of n and delta equals the probability of rejecting the null hypothesis given delta and n"

\textbf{What it does:} Calculates the statistical power for a given sample size and effect size

\textbf{Why peppwR uses it:} This answers "What power will my planned experiment have?"

\subsection{Formula 5: Minimum Detectable Effect Size}

\textbf{Mathematical notation:}
\[\delta^* = \min\{\delta : P(\text{reject } H_0 | \delta, n) \geq \pi_{\text{target}}\}\]

\textbf{What the symbols mean:}
\begin{itemize}
\item $\delta^*$ = minimum detectable effect size
\item Other symbols as defined previously
\end{itemize}

\textbf{How to read it:} "Delta-star equals the minimum delta such that the probability of rejecting the null hypothesis, given delta and n, is greater than or equal to target power"

\textbf{What it does:} Finds the smallest biological effect you can reliably detect with your sample size

\textbf{Why peppwR uses it:} This answers "What's the smallest effect I can detect?" given your experimental constraints

\subsection{Formula 6: Peptidome-wide Power}

\textbf{Mathematical notation:}
\[P_{\text{peptidome}}(n, \delta, \pi_{\text{target}}) = \frac{1}{K}\sum_{i=1}^{K} \mathbf{1}[\pi_i(n, \delta) \geq \pi_{\text{target}}]\]

\textbf{What the symbols mean:}
\begin{itemize}
\item $P_{\text{peptidome}}$ = proportion of peptides achieving target power
\item $K$ = total number of peptides
\item $\sum_{i=1}^{K}$ = sum from peptide 1 to peptide K
\item $\mathbf{1}[\cdot]$ = indicator function (equals 1 if condition is true, 0 if false)
\item $\pi_i(n, \delta)$ = power for peptide i
\end{itemize}

\textbf{How to read it:} "Peptidome power equals one over K times the sum of indicator functions for whether each peptide's power exceeds the target"

\textbf{What it does:} Calculates the fraction of peptides that will achieve your target power level

\textbf{Why peppwR uses it:} Since different peptides have different properties, this tells you what percentage of your peptides will be detectable

\subsection{Formula 7: FDR Control (Benjamini-Hochberg)}

\textbf{Mathematical notation:}
\[k^* = \max\left\{i : p_{(i)} \leq \frac{i \cdot \alpha}{m}\right\}\]

\textbf{What the symbols mean:}
\begin{itemize}
\item $k^*$ = cutoff index for rejecting hypotheses
\item $\max\{\cdot\}$ = maximum value satisfying the condition
\item $p_{(i)}$ = i-th smallest p-value (ordered from smallest to largest)
\item $\alpha$ = significance level (usually 0.05)
\item $m$ = total number of statistical tests performed
\end{itemize}

\textbf{How to read it:} "k-star equals the maximum i such that the i-th ordered p-value is less than or equal to i times alpha divided by m"

\textbf{What it does:} Determines which p-values to call significant when testing many peptides simultaneously

\textbf{Why peppwR uses it:} When testing thousands of peptides, you need to adjust for multiple comparisons to avoid false discoveries

\section{Distribution Fitting Section}

\subsection{Formula 8: Maximum Likelihood Estimation}

\textbf{Mathematical notation:}
\[\hat{\theta}_{k,i} = \arg\max_{\theta_k} \sum_{j=1}^{n_i} \log f_k(x_{ij} | \theta_k)\]

\textbf{What the symbols mean:}
\begin{itemize}
\item $\hat{\theta}_{k,i}$ = estimated parameters for distribution k and peptide i
\item $\arg\max$ = "the argument that maximizes" - find the parameter values that make the expression biggest
\item $\theta_k$ = parameters of distribution k
\item $\sum_{j=1}^{n_i}$ = sum over all j from 1 to $n_i$ (all observations for peptide i)
\item $\log$ = natural logarithm
\item $f_k(x_{ij} | \theta_k)$ = probability density function for distribution k
\item $x_{ij}$ = j-th observation for peptide i
\end{itemize}

\textbf{How to read it:} "Theta-hat k,i equals the argument that maximizes the sum of log probability densities"

\textbf{What it does:} Finds the distribution parameters that make the observed data most likely

\textbf{Why peppwR uses it:} This is how we fit mathematical distributions (like gamma or lognormal) to each peptide's abundance data

\subsection{Formula 9: Model Selection (AIC)}

\textbf{Mathematical notation:}
\[\hat{k}_i = \arg\min_k \left\{-2 \sum_{j=1}^{n_i} \log f_k(x_{ij} | \hat{\theta}_{k,i}) + 2p_k\right\}\]

\textbf{What the symbols mean:}
\begin{itemize}
\item $\hat{k}_i$ = best distribution choice for peptide i
\item $\arg\min$ = "the argument that minimizes" - find the k that makes the expression smallest
\item $-2$ = constant multiplier in AIC formula
\item $2p_k$ = penalty term where $p_k$ is the number of parameters in distribution k
\end{itemize}

\textbf{How to read it:} "k-hat i equals the argument that minimizes negative two times the log-likelihood plus two times the number of parameters"

\textbf{What it does:} Chooses the best distribution for each peptide, balancing goodness of fit against complexity

\textbf{Why peppwR uses it:} AIC helps select whether gamma, lognormal, normal, or inverse Gaussian best describes each peptide's data

\section{Simulation Engine Section}

\subsection{Formula 10: Monte Carlo Power Estimation}

\textbf{Mathematical notation:}
\[\hat{\pi}_i(n, \delta) = \frac{1}{n_{\text{sim}}} \sum_{j=1}^{n_{\text{sim}}} R^{(j)}\]

\textbf{What the symbols mean:}
\begin{itemize}
\item $\hat{\pi}_i(n, \delta)$ = estimated power for peptide i
\item $n_{\text{sim}}$ = number of simulations (usually 1000)
\item $R^{(j)}$ = result of simulation j (1 if significant, 0 if not)
\item $\sum_{j=1}^{n_{\text{sim}}}$ = sum over all simulations
\end{itemize}

\textbf{How to read it:} "Pi-hat i equals one over n-sim times the sum of all simulation results"

\textbf{What it does:} Estimates power by running many simulated experiments and counting how often they detect the effect

\textbf{Why peppwR uses it:} Since analytical power calculations are impossible for complex proteomics data, simulation provides accurate estimates

\subsection{Formula 11: Monte Carlo Standard Error}

\textbf{Mathematical notation:}
\[\text{SE}(\hat{\pi}_i) = \sqrt{\frac{\hat{\pi}_i(1-\hat{\pi}_i)}{n_{\text{sim}}}}\]

\textbf{What the symbols mean:}
\begin{itemize}
\item $\text{SE}(\hat{\pi}_i)$ = standard error of the power estimate
\item $\sqrt{\cdot}$ = square root
\item $(1-\hat{\pi}_i)$ = one minus the power estimate
\end{itemize}

\textbf{How to read it:} "Standard error of pi-hat equals the square root of pi-hat times one minus pi-hat, divided by the number of simulations"

\textbf{What it does:} Calculates the uncertainty in our power estimate due to simulation randomness

\textbf{Why peppwR uses it:} This tells you how precise your power estimate is - more simulations give lower standard error

\section{Missing Data Handling Section}

\subsection{Formula 12: Mean Abundance Calculation}

\textbf{Mathematical notation:}
\[\bar{x}_i = \frac{1}{n_{\text{obs},i}} \sum_{j: x_{ij} \neq \text{NA}} x_{ij}\]

\textbf{What the symbols mean:}
\begin{itemize}
\item $\bar{x}_i$ = mean abundance for peptide i
\item $n_{\text{obs},i}$ = number of observed (non-missing) values for peptide i
\item $\sum_{j: x_{ij} \neq \text{NA}}$ = sum over only the non-missing observations
\item $\text{NA}$ = missing value (Not Available)
\end{itemize}

\textbf{How to read it:} "x-bar i equals one over n-observed,i times the sum of all non-missing x values"

\textbf{What it does:} Calculates the average abundance for a peptide, excluding missing values

\textbf{Why peppwR uses it:} Missing data is common in proteomics, so we calculate means from available data only

\subsection{Formula 13: Missingness Rate}

\textbf{Mathematical notation:}
\[r_i = \frac{n_{\text{missing},i}}{n_{\text{total},i}}\]

\textbf{What the symbols mean:}
\begin{itemize}
\item $r_i$ = missingness rate for peptide i
\item $n_{\text{missing},i}$ = number of missing values for peptide i
\item $n_{\text{total},i}$ = total number of samples for peptide i
\end{itemize}

\textbf{How to read it:} "r i equals the number of missing values divided by the total number of values"

\textbf{What it does:} Calculates what fraction of data is missing for each peptide

\textbf{Why peppwR uses it:} Different peptides have different amounts of missing data, which affects power calculations

\subsection{Formula 14: Missing Data in Simulations}

\textbf{Mathematical notation:}
\[\lfloor n \cdot r_i \rfloor\]

\textbf{What the symbols mean:}
\begin{itemize}
\item $\lfloor \cdot \rfloor$ = floor function (rounds down to nearest integer)
\item $n$ = sample size
\item $r_i$ = missingness rate for peptide i
\end{itemize}

\textbf{How to read it:} "Floor of n times r i"

\textbf{What it does:} Calculates how many observations to randomly remove from simulated data

\textbf{Why peppwR uses it:} This recreates realistic missing data patterns in power simulations

\section{Summary}

This guide covers all major mathematical notation in the peppwR Methods section. The formulas fall into four main categories:

\begin{enumerate}
\item \textbf{Statistical Framework:} Defines power, hypothesis tests, and optimization problems
\item \textbf{Distribution Fitting:} Uses maximum likelihood and AIC for model selection
\item \textbf{Simulation Engine:} Estimates power through Monte Carlo simulation
\item \textbf{Missing Data:} Characterizes and incorporates missingness patterns
\end{enumerate}

The mathematical notation enables precise description of peppwR's algorithms while the simulations handle the complexity that makes analytical solutions impossible for proteomics data.

\end{document}